\chapter{Technology Strategy}\label{ch:fm:technology-strategy}

Virtu Financial's communication and data processing expense was \$209 million in 2019 and \$249 million in 2025. In between, its revenue went from
\$1.5 billion to \$3.2 billion, back down to \$2.3 billion and up to \$3.6 billion; the expense rose by a fifth and fell only once, by less than 1\%. A trading firm does not
choose each year whether to be fast. The \omterm{def:fm:technology-strategy:tier}{latency tier} it pays for is a decision about which businesses it will be in, taken once and paid every
year, whatever the year brings.

\section{Architecture as strategy}

Book 13 built the low-latency stack and Book 14 the networks under it. This chapter is about the decision above both: how much speed to buy, for
which strategies, and what else the same money could do. The decision is strategic because the costs are fixed (chapter 1's \omterm{def:fm:the-economics-of-a-trading-firm:fixed}{fixed cost base}; \cref{fig:fm:technology-strategy:virtu}) and
the benefits depend on what competitors buy.

\begin{omfigure}
\begin{tikzpicture}
\begin{axis}[width=10.5cm,height=5.2cm,ymode=log,xlabel={\small year},ylabel={\small \$ million (log scale)},xmin=2018.6,xmax=2025.4,ymin=150,
  ymax=5000,xtick={2019,2020,2021,2022,2023,2024,2025},xticklabel style={/pgf/number format/1000 sep=},ytick={200,500,1000,2000,4000},
  yticklabels={200,500,1000,2000,4000},tick label style={font=\footnotesize},grid=major,grid style={gray!20},table/col sep=comma,
  legend cell align=left,legend style={font=\scriptsize,at={(0.5,-0.3)},anchor=north,/tikz/every even column/.append style={column sep=8pt}},legend columns=2]
\addplot[omDef,thick,mark=*,mark size=1.3pt] table[x=year,y=revenue]{figdata/desk/19-technology-strategy/virtu.csv};
\addplot[omProp,thick,mark=square*,mark size=1.3pt] table[x=year,y=comm_data]{figdata/desk/19-technology-strategy/virtu.csv};
\legend{total revenue,communication and data processing}
\end{axis}
\end{tikzpicture}

\omcaption{Virtu Financial's total revenue and its communication and data processing expense, 2019--2025. Revenue more than doubled and fell back;
the expense rose steadily. Source: Forms 10-K (the chapter's ledger); data in \texttt{data/desk/filings\_virtu.csv}.}\label{fig:fm:technology-strategy:virtu}
\end{omfigure}

\begin{definition}[Latency tier]\label{def:fm:technology-strategy:tier}
A \emph{latency tier}\index{latency tier} is a level of trading infrastructure defined by the latency it achieves and the fixed cost it carries:
for example, shared connections and software trading; colocated servers with direct low-latency connections; or colocated hardware trading with
dedicated wireless routes between venues.
\end{definition}

The chapter's three tiers are priced from venue fees Book 14 read from published fee schedules, plus the firm's own hardware, routes and people,
which are inputs (\cref{dat:fm:technology-strategy:fees}): \$0.29 million a year for tier 1, \$1.35 million for tier 2 and \$4.45 million for tier 3,
per strategy. The venue fees are the small part: \$36\,000, \$450\,000 and \$946\,000 a year.

\begin{dated}{2026-09}{Venue fees behind the tiers}\label{dat:fm:technology-strategy:fees}
From the fee schedules recorded in One Quant Book 14 (chapter 9): at Nasdaq's NY11-4 data centre, an ultra-high-density cabinet at \$7\,230 a month,
with installation fees of \$5\,940 for the cabinet, \$4\,560 for Phase 3 power and \$5\,260 for a Phase 3 PDU (SEC Release 34-101267, October
2024); at MIAX Pearl Options, a 10Gb ultra-low-latency connection at \$15\,000 a month, a 10Gb disaster-recovery connection at \$3\,500 and a 1Gb
connection at \$1\,500 (Federal Register 2026-10452, May 2026). Tier 1 is two 1Gb connections; tier 2 a cabinet and two 10Gb connections; tier 3
two cabinets, four 10Gb connections and a disaster-recovery link, one-time fees spread over 36 months. Fee schedules change by filing.
\end{dated}

\section{What speed is worth, strategy by strategy}

Speed earns money in races: two or more firms react to the same public event, and only the first gets the trade. Aquilina, Budish and O'Neill
(2022) measured races directly in stock exchange message data for FTSE 100 stocks: about one race a minute per symbol, lasting 5 to 10
millionths of a second in the modal case, about a fifth of trading volume, with six firms accounting for over 80\% of race wins and losses. Races
were roughly a third of the cost of liquidity. For a strategy, the question is what share of its revenue is decided by races.

\begin{definition}[Latency arms race]\label{def:fm:technology-strategy:race}
A \emph{latency arms race}\index{latency arms race} is competition in which firms buy speed to win races against each other, so that each firm's
investment lowers the others' returns on theirs, and the equilibrium spending on speed exceeds what the firms together gain from it.
\end{definition}

The chapter's race model splits a strategy's revenue pool $V$, shared by $n$ competitors, into a race share $\alpha$ and the rest. Races go to the
firms in the fastest tier present, shared equally; the rest is shared equally among all $n$. A firm earns
\[
V\Bigl(\alpha\,\frac{\mathbf 1\{\text{in the top tier}\}}{m}+\frac{1-\alpha}{n}\Bigr)-c_{\text{tier}},
\]
where $m$ is the number of firms in the top tier (\cref{lst:fm:technology-strategy:race}). Five strategies differ in $V$ and in $\alpha$: futures
market making (\$60 million, 0.8), ETF arbitrage (\$30 million, 0.9), equity statistical arbitrage over minutes (\$20 million, 0.1), options
market making (\$40 million, 0.5) and crypto basis trading (\$10 million, 0.3). The values are illustrative; the race shares are set against the
finding that races are a fifth of volume and a third of the cost of liquidity in a market where market makers dominate both.

\omcode{code/firm/techtier/firm_techtier.py}{52}{76}{The race model: a firm's revenue given its tier and its rivals', its best tier, and the arms
race solved by best responses.}\label{lst:fm:technology-strategy:race}

\section{Tutorial: the fastest tier}\label{tut:fm:technology-strategy:tut}

\begin{tutorial}
\textbf{Goal.} Choose a \omterm{def:fm:technology-strategy:tier}{latency tier} for each of five strategies against given rivals, then let every firm choose and find the arms-race
equilibrium. \textbf{End state:} the profit chart by tier (\cref{fig:fm:technology-strategy:profits}) and equilibrium spending against no upgrades
(\cref{fig:fm:technology-strategy:race}).
\begin{tutsteps}
\item \textbf{The tiers.} \texttt{fm\_tech.tiers()} builds the three tiers from Book 14's \texttt{firm.colobill} schedules and the stated inputs.
\item \textbf{The rivals.} Each strategy has four rivals at given tiers: for futures market making two at tier 3, one at tier 2 and one at tier 1.
\item \textbf{The choice.} \texttt{firm.techtier.best\_tier} compares the three tiers' profits for each strategy.
\item \textbf{The race.} \texttt{firm.techtier.arms\_race} starts five identical firms at tier 1 and lets each in turn choose its best tier until no
one changes; \texttt{waste} compares the equilibrium's spending with everyone at tier 1.
\end{tutsteps}
\end{tutorial}

\begin{omfigure}
\begin{tikzpicture}
\begin{axis}[width=10.5cm,height=5.4cm,ybar,bar width=7pt,ylabel={\small profit (\$ million a year)},ymin=-2,ymax=22,xtick={0,1,2,3,4},
  xticklabels={{futures MM},{ETF arb},{stat arb},{options MM},{crypto}},xmin=-0.6,xmax=4.6,
  tick label style={font=\footnotesize},grid=major,grid style={gray!20},table/col sep=comma,
  legend cell align=left,legend style={font=\scriptsize,at={(0.5,-0.25)},anchor=north,/tikz/every even column/.append style={column sep=8pt}},legend columns=3]
\addplot[fill=gray!40,draw=gray] table[x=pos,y=t1]{figdata/desk/19-technology-strategy/profits.csv};
\addplot[fill=omThm!50,draw=omThm] table[x=pos,y=t2]{figdata/desk/19-technology-strategy/profits.csv};
\addplot[fill=omDef!55,draw=omDef] table[x=pos,y=t3]{figdata/desk/19-technology-strategy/profits.csv};
\legend{tier 1,tier 2,tier 3}
\end{axis}
\end{tikzpicture}

\omcaption{Profit by \omterm{def:fm:technology-strategy:tier}{latency tier} for each of five strategies, against given rivals. Speed pays where races decide most of the revenue and the pool
is large; tier 3 earns less for the statistical arbitrage book and loses money on the crypto basis trade. Data: \texttt{fm\_tech.choices}.}\label{fig:fm:technology-strategy:profits}
\end{omfigure}

Against its rivals, the firm should buy tier 3 for futures market making (profit \$13.95 million a year, against \$2.11 million at tier 1), ETF
arbitrage (\$9.65 million against \$0.31 million) and options market making (\$19.55 million against \$3.71 million); tier 2 for the crypto basis
trade (\$3.05 million), where no rival has more; and tier 1 for statistical arbitrage over minutes (\$3.31 million), which loses \$2.16 million a
year by buying tier 3. The same firm is a speed business in three strategies and not in two, and its architecture should say so.

\begin{proposition}[How many firms race]\label{prop:fm:technology-strategy:count}
With two tiers whose costs differ by $\Delta c$, $n$ identical firms and the race model, a firm gains from joining $m-1$ others in the fast tier if
and only if $\alpha V/m\ge\Delta c$. In equilibrium the number of firms in the fast tier is $m^*=\min(n,\lfloor\alpha V/\Delta c\rfloor)$, their
total speed spending is $m^*\Delta c$ above the all-slow level, and the pool shared by the industry is unchanged.
\end{proposition}

\begin{proof}
Moving from the slow to the fast tier changes a firm's race revenue from 0 to $\alpha V/m$ and leaves its share of the rest at $(1-\alpha)V/n$;
the cost rises by $\Delta c$. Firms join while $\alpha V/m\ge\Delta c$, so the last to join is the $\lfloor\alpha V/\Delta c\rfloor$-th, and a firm
in the fast tier does not leave, since leaving loses $\alpha V/m^*\ge\Delta c$. Revenues sum to $V$ whatever the tiers.
\end{proof}

With three tiers the argument applies step by step: for futures market making $\alpha V=48$ and the step from tier 1 to tier 3 is \$4.16 million,
so up to eleven firms would race; with five, all five do. Options market making gives $20/4.16$, four firms; the crypto trade, $3/1.06$ for the
step to tier 2, two firms; statistical arbitrage, $2/1.06$, one.

\begin{omfigure}
\begin{tikzpicture}
\begin{axis}[width=10.5cm,height=5.2cm,ybar,bar width=10pt,ylabel={\small speed spending (\$ million a year)},ymin=0,ymax=25,xtick={0,1,2,3,4},
  xticklabels={{futures MM},{ETF arb},{stat arb},{options MM},{crypto}},xmin=-0.6,xmax=4.6,
  tick label style={font=\footnotesize},grid=major,grid style={gray!20},table/col sep=comma,
  legend cell align=left,legend style={font=\scriptsize,at={(0.5,-0.25)},anchor=north,/tikz/every even column/.append style={column sep=8pt}},legend columns=2]
\addplot[fill=omDef!55,draw=omDef] table[x=pos,y=spend]{figdata/desk/19-technology-strategy/race.csv};
\addplot[fill=gray!40,draw=gray] table[x=pos,y=baseline]{figdata/desk/19-technology-strategy/race.csv};
\legend{arms-race equilibrium,no one upgrades}
\end{axis}
\end{tikzpicture}

\omcaption{Industry spending on speed by five identical firms in each strategy, at the arms-race equilibrium and if no one upgraded. The revenue
pool is the same in both; the difference is spent to redistribute it. Data: \texttt{fm\_tech.races}.}\label{fig:fm:technology-strategy:race}
\end{omfigure}

\section{Latency tiers and the arms race}

The equilibrium spends \$22.2 million a year on speed in futures market making and in ETF arbitrage, against \$1.4 million if no one upgraded: 93.6\%
of the spending buys nothing for the industry, since the pool is the same. The industry's profit in futures market making falls from \$58.6 million
to \$37.8 million; in ETF arbitrage, from \$28.6 million to \$7.8 million, about a quarter of what it would be. Options market making wastes 92.1\%, the
crypto trade 59.8\% and statistical arbitrage 42.7\%. Across the five, the industry spends \$68.6 million where \$7.2 million would do: 89.6\% of
its speed spending is waste in the model's sense. This is the argument Budish, Cramton and Shim (2015) make for markets in general: competition does
not shrink the arbitrage opportunities, it raises the speed needed to capture them.

For a single firm the conclusion is the opposite of the industry's: in a race that others are running, not buying speed means losing the race's
revenue. The firm's decision is where to race and where not to, and where the pool is small or the races few, to leave.

\section{Platform teams and product teams}

\begin{definition}[Platform team, product team]\label{def:fm:technology-strategy:teams}
A \emph{product team}\index{product team} builds and runs what one business uses to make money, such as a strategy's trading system, and is
measured on that business's results. A \emph{platform team}\index{platform team} builds and runs shared capabilities that product teams use as a
service, such as market data, order gateways, the research platform or deployment, and is measured on their reliability and on the product teams'
speed of delivery.
\end{definition}

The \omterm{def:fm:technology-strategy:tier}{latency tier} cuts across the two. Tier 3 infrastructure is usually a platform, since venues, cabinets and routes are shared by every strategy
that races; the tutorial charges each strategy its own tier only for clarity. What belongs in a platform is what several \omterm{def:fm:technology-strategy:teams}{product teams} need in the
same form; what belongs in a \omterm{def:fm:technology-strategy:teams}{product team} is what makes a strategy different. Skelton and Pais's \emph{Team Topologies} (2019) gives the
vocabulary; chapter 21 applies it to a trading firm's engineers (\cref{fig:fm:technology-strategy:teams}).

\begin{omfigure}
\begin{tikzpicture}[font=\footnotesize,>=stealth,box/.style={draw,rounded corners,align=center,minimum height=0.75cm,text width=2.3cm}]
\foreach \x/\n in {-4.2/futures MM,-1.4/ETF arb,1.4/options MM,4.2/stat arb} {\node[box,fill=omDef!15] at (\x,1.5) {\n\\{\scriptsize product team}};}
\node[box,fill=omThm!15,text width=7.6cm] (lat) at (-1.4,0) {low-latency platform: venues, cabinets, routes, gateways (tier 3)};
\node[box,fill=gray!15,text width=10.6cm] (base) at (0,-1.1) {base platform: market data, deployment, research, monitoring};
\draw[->] (-4.2,0.4) -- (-4.2,1.1); \draw[->] (-1.4,0.4) -- (-1.4,1.1); \draw[->] (1.4,0.4) -- (1.4,1.1);
\draw[->] (4.2,-0.7) -- (4.2,1.1);
\end{tikzpicture}

\omcaption{Platform and \omterm{def:fm:technology-strategy:teams}{product teams} for the chapter's firm: three racing strategies share a low-latency platform; the statistical arbitrage team
uses only the base platform. Schematic.}\label{fig:fm:technology-strategy:teams}
\end{omfigure}

\section{The technology budget}

\begin{definition}[Run-the-bank and change-the-bank spending]\label{def:fm:technology-strategy:runchange}
\emph{Run-the-bank spending}\index{run-the-bank spending} is the technology spending needed to keep existing businesses operating: infrastructure,
licences, venue and data fees, production support and maintenance. \emph{Change-the-bank spending}\index{change-the-bank spending} is spending on
new capabilities: new strategies, venues, platforms and replacements that change what the firm can do.
\end{definition}

The chosen tiers cost \$14.97 million a year and production support \$1.2 million: \$16.2 million of run. The year's projects, new venue
connections, an FPGA feed-handler upgrade and a research platform, cost \$4.0 million of change: run is 80.2\% of the budget. The hook's filings
show why the split matters: a cost that did not fall when revenue fell by 29\% from 2020 to 2023 is a run cost, and every tier decision adds to
it. Flow Traders' annual report puts its technology expenses at \euro70.6 million in 2025, against \euro66.6 million in 2024, beside a net trading
income of \euro485.8 million: another fixed line in a variable business.

\begin{method}[Setting the technology strategy]\label{met:fm:technology-strategy:set}
\begin{enumerate}
\item For each strategy, estimate the share of revenue decided by races and the competitors' tiers.
\item Price each tier from published fees and the firm's own costs; choose the tier that maximises the strategy's profit, and record the choice
as a strategic commitment with a review date.
\item Share racing infrastructure as a platform; keep strategy-specific code in \omterm{def:fm:technology-strategy:teams}{product teams}.
\item Report the budget as run and change, and require every change project to state the run cost it will add.
\item Revisit the choice when competitors move: the arms race changes the answer every year.
\end{enumerate}
\end{method}

\section{Build: latency tiers}\label{bld:fm:technology-strategy:techtier}

\begin{build}
\bfield{Purpose} The strategic side of speed: tier costs, what each tier earns strategy by strategy, the arms-race equilibrium, and the run/change
budget.
\bfield{Interface} \texttt{firm.techtier}: \texttt{Tier}, \texttt{colobill\_venue\_fees} (on Book 14's \texttt{firm.colobill}), \texttt{Strategy},
\texttt{revenue}, \texttt{best\_tier}, \texttt{arms\_race}, \texttt{waste}, \texttt{run\_change}.
\bfield{Rules} Revenues sum to the pool; the fastest tier present takes all races, shared equally; an equilibrium is a profile no firm changes;
cycling best responses raise an error rather than return a profile.
\bfield{Acceptance tests} \texttt{code/firm/techtier/tests/}: the pool identity; no profitable deviation at the equilibrium; venue fees from a
published schedule; the budget split.
\bfield{Stretch} Book 14's connectivity plan as the source of tier costs; shared platform costs across strategies; mixed equilibria when best
responses cycle.
\end{build}

\begin{omsources}
\item E.~Budish, P.~Cramton and J.~Shim, ``The high-frequency trading arms race'', \emph{Quarterly Journal of Economics} 130(4), 2015;
M.~Aquilina, E.~Budish and P.~O'Neill, ``Quantifying the high-frequency trading `arms race''', \emph{Quarterly Journal of Economics} 137(1), 2022.
\item Virtu Financial, Forms 10-K 2020--2025; Flow Traders, Annual Report 2025.
\item M.~Skelton and M.~Pais, \emph{Team Topologies}, IT Revolution, 2019.
\item Fee schedules as recorded in One Quant Book 14, chapter 9.
\end{omsources}

\section{Exercises}

\begin{exercise}[$\star$]\label{exo:fm:technology-strategy:1}
By what percentage did Virtu's communication and data expense rise from 2019 to 2025, and what share of revenue was it in each year?
\end{exercise}

\begin{exercise}[$\star$]\label{exo:fm:technology-strategy:2}
A strategy has $V=60$, $\alpha=0.8$ and five firms. What does a firm earn in the fast tier alone, and with two others there?
\end{exercise}

\begin{exercise}[$\star$]\label{exo:fm:technology-strategy:3}
What are the venue fees of the three tiers, and what share of each tier's cost are they?
\end{exercise}

\begin{exercise}[$\star\star$]\label{exo:fm:technology-strategy:4}
Use \cref{prop:fm:technology-strategy:count} to predict how many firms race in each strategy.
\end{exercise}

\begin{exercise}[$\star\star$]\label{exo:fm:technology-strategy:5}
Why should the firm not buy tier 3 for statistical arbitrage although it buys it for three other strategies?
\end{exercise}

\begin{exercise}[$\star\star$]\label{exo:fm:technology-strategy:6}
Classify as run or change: renewing a cabinet contract, building a new venue gateway, replacing an end-of-life switch, the market data bill.
\end{exercise}

\begin{exercise}[$\star\star\star$]\label{exo:fm:technology-strategy:7}
\emph{Coding.} Double tier 3's other costs to \$7 million and rerun \texttt{fm\_tech.races()}. What happens to the equilibria and the waste?
\end{exercise}

\begin{exercise}[$\star\star\star$]\label{exo:fm:technology-strategy:8}
\emph{Find the flaw.} ``Speed spending is wasted, as the models show; we will stop upgrading and save the money.''
\end{exercise}

\section{Problem: The Fastest Tier}

\begin{problem}[{Weekend problem --- the fastest tier}]\label{pb:fm:technology-strategy:1}
A chief technology officer must propose which strategies the firm will race in, and at what speed, for the next three years.

\medskip\noindent\textbf{Part I --- Speed as strategy.}
\begin{enumerate}
\item Define a \omterm{def:fm:technology-strategy:tier}{latency tier} and a \omterm{def:fm:technology-strategy:race}{latency arms race}.
\item What does Virtu's communication and data expense show?
\item What did Aquilina, Budish and O'Neill find about races?
\item Price the three tiers and state where the prices come from.
\end{enumerate}

\medskip\noindent\textbf{Part II --- The choice.}
\begin{enumerate}[resume]
\item Write the race model.
\item Give the profit by tier for each strategy and the best tier.
\item Why is the answer different for the crypto basis trade?
\item What would change the answer for statistical arbitrage?
\end{enumerate}

\medskip\noindent\textbf{Part III --- The race.}
\begin{enumerate}[resume]
\item State and prove \cref{prop:fm:technology-strategy:count}.
\item Give each strategy's equilibrium profile and spending.
\item Give the wasted share per strategy and in total.
\item Why is the industry's waste not the firm's?
\item What do Budish, Cramton and Shim conclude?
\end{enumerate}

\medskip\noindent\textbf{Part IV --- The organisation and the budget.}
\begin{enumerate}[resume]
\item Define platform and \omterm{def:fm:technology-strategy:teams}{product teams} and draw them for the firm.
\item Define run-the-bank and \omterm{def:fm:technology-strategy:runchange}{change-the-bank spending}.
\item Give the firm's run/change split.
\item What does Flow Traders' report show?
\item How often should the tier choice be revisited, and on what evidence?
\item State the \emph{named result}: the profit-maximising tier for each of the five strategies, and the share of the industry's speed spending that
the arms-race equilibrium wastes against no one upgrading.
\item In two sentences, write the technology strategy.
\end{enumerate}
\end{problem}

\section{Interview questions}

\begin{interviewq}[$\star$ \iqroles{developer}]\label{iq:fm:technology-strategy:1}
What is the difference between a \omterm{def:fm:technology-strategy:teams}{platform team} and a \omterm{def:fm:technology-strategy:teams}{product team} in a trading firm?
\end{interviewq}

\begin{interviewq}[$\star$ \iqroles{trader}]\label{iq:fm:technology-strategy:2}
Why can a firm lose money by being faster?
\end{interviewq}

\begin{interviewq}[$\star\star$ \iqroles{researcher}]\label{iq:fm:technology-strategy:3}
Five firms compete for a \$30 million pool, 90\% decided by races; being fast costs \$4 million more a year. How many will be fast?
\end{interviewq}

\begin{interviewq}[$\star\star$ \iqroles{developer, trader}]\label{iq:fm:technology-strategy:4}
How would you decide whether a strategy needs an FPGA?
\end{interviewq}

\begin{interviewq}[$\star\star$ \iqroles{risk}]\label{iq:fm:technology-strategy:5}
Why does most of a trading firm's technology budget become \omterm{def:fm:technology-strategy:runchange}{run-the-bank spending}?
\end{interviewq}

\begin{interviewq}[$\star\star\star$ \iqroles{researcher}]\label{iq:fm:technology-strategy:6}
Market designers propose frequent batch auctions. What would they do to the arms race in the chapter's model?
\end{interviewq}
